% Chapter 1 -- Rubric: PO1 (1.1), PO2 (problem identification), PO3 (3.2 objectives)
\chapter{Problem Identification and Formulation}
\label{ch:problem}

\section{Background and Motivation}
\label{sec:background}

Admission to the engineering universities of Bangladesh --- BUET, KUET, RUET
and CUET --- is decided by a single competitive examination taken after the
Higher Secondary Certificate (HSC). The test covers Mathematics, Physics and
Chemistry, and the ratio of candidates to available seats is severe enough that
preparation is effectively a year-long full-time activity for those who attempt
it. The dominant preparation method is volume: a student buys one or more
question banks compiled from previous years' papers and works through them,
usually alongside coaching-centre classes that move at a fixed pace through a
fixed syllabus.

This method has a structural weakness that has nothing to do with effort. When
a student fails a question, the question bank tells them the correct answer; it
does not tell them \emph{why} they failed. A wrong answer on a problem
requiring integration by parts may indicate that the student does not recognise
when the technique applies, or cannot choose $u$ and $dv$ correctly, or chose
them correctly but made an algebraic error in the resulting simpler integral,
or does not know the derivative of the function involved. These are four
different deficits requiring four different remedies, and they are invisible in
a score. A human tutor localises them by watching the student work --- which is
precisely what one-to-one tutoring is effective at, and precisely what does not
scale to a national cohort \cite{vanlehn2011relative}.

The consequence is misallocated study time. A student who does not know which
of their foundations is weak has no rational basis for choosing what to study
next, and defaults to either linear syllabus coverage or to practising the
topics where failure is most visible --- often the topics \emph{downstream} of
the actual gap. Time spent re-attempting integration problems does not repair
an algebra deficit; it merely produces repeated failure on the same
prerequisite, which is both inefficient and demoralising.

What is needed is a system that infers, from the student's own answers, where
in the structure of the subject their knowledge stops --- and then selects the
next question at that frontier rather than at the point in the syllabus the
calendar happens to have reached.

\section{Problem Statement}
\label{sec:problem-statement}

\textbf{Problem.} Given a learner and a sequence of their responses to
multiple-choice questions, estimate the learner's mastery of every skill in a
subject domain, and select the next question so as to localise and repair the
learner's knowledge frontier.

Formulated in engineering terms:

\begin{itemize}
  \item \textbf{Entities.} A set of \emph{skills} $S = \{s_1,\dots,s_n\}$, each
    an atomic, independently testable unit of knowledge or procedure, annotated
    with one Bloom's Taxonomy cognitive level. A set of directed
    \emph{prerequisite} edges $E \subseteq S \times S$, where $(s_i,s_j) \in E$
    means $s_i$ must be held before $s_j$ is reachable. A bank of questions,
    each tagged with the skill it tests, its Bloom level and its topic. A set of
    learners.
  \item \textbf{Inputs.} For a learner $u$, an observed sequence of binary
    outcomes $(q_1,c_1),(q_2,c_2),\dots$ where $q_t$ is the question served at
    step $t$ and $c_t \in \{\text{correct},\text{incorrect}\}$.
  \item \textbf{Latent state.} For each skill, $P(L_{u,s}) \in [0,1]$, the
    probability that learner $u$ has learned skill $s$. This is never observed
    directly; it must be inferred from noisy evidence.
  \item \textbf{Outputs.} (i) The posterior $P(L_{u,s})$ for all $s \in S$
    after each response; (ii) the identity of the next question to serve.
  \item \textbf{Quantities to be guaranteed.} The graph $(S,E)$ must be acyclic,
    or the notion of ``prerequisite'' is incoherent and propagation does not
    terminate. Mastery estimates must respect the graph: no skill may be
    credited as mastered while an ancestor it depends on is estimated as
    unmastered. Question selection must terminate even when the bank has no
    item matching the ideal specification.
  \item \textbf{Assumptions.} Responses are conditionally independent given the
    latent state; a learner does not forget within a session (no decay term);
    each question tests primarily one skill; the prerequisite structure is the
    same for all learners.
\end{itemize}

\noindent\textbf{Engineering knowledge demanded.} The formulation above is not
solvable by routine application of standard components. It requires: probability
theory and Bayesian inference (the latent-state update); graph theory (acyclicity
checking, topological ordering, transitive reduction, ancestor traversal);
algorithm design (the question-selection heuristic and its fallback ladder);
database design under referential constraints; the pedagogical literature on
knowledge tracing and cognitive taxonomies, which supplies the model family and
the level annotation scheme; and, for the content pipeline, applied natural
language processing --- retrieval-augmented generation, embedding-based
retrieval, and the validation of machine-generated assessment items.

\section{Complexity of the Problem}
\label{sec:complexity}
A complex engineering problem must have characteristic P1 and some or all of
P2--P7 (\Cref{tab:complexity}). P1 requires in-depth knowledge at the level
of one or more of the knowledge profiles K3--K8 listed in
\Cref{tab:knowledge-profile}.

% Rendered as a longtable (not a float) so it flows across pages at readable
% size instead of being squeezed onto one rotated page.
\setlength{\LTcapwidth}{\textwidth}
{\small\setstretch{1.1}
\begin{longtable}{@{}p{11mm}p{56mm}p{\dimexpr\textwidth-11mm-56mm-4\tabcolsep\relax}@{}}
  \caption{Complex engineering problem characteristics (P1--P7) and evidence from this project}
  \label{tab:complexity}\\
    \toprule
    Code & Definition & Evidence from this project \\
    \midrule
  \endfirsthead
    \multicolumn{3}{@{}l}{\small\itshape \tablename~\thetable{} --- continued from previous page}\\[2pt]
    \toprule
    Code & Definition & Evidence from this project \\
    \midrule
  \endhead
    \midrule
    \multicolumn{3}{r@{}}{\small\itshape continued on next page}\\
  \endfoot
    \bottomrule
  \endlastfoot
    P1 & \textbf{Depth of knowledge required.} Cannot be resolved without in-depth engineering knowledge at the level of one or more of K3, K4, K5, K6, or K8, which allows a fundamentals-based, first-principles analytical approach. &
    \textbf{K3, K4, K5, K8.} K3: the Bayesian update is derived from first principles, not called from a library --- the posterior and transition equations are implemented directly (\Cref{sec:bkt-model}). K4: knowledge tracing and cognitive-taxonomy literature supplies the model family \cite{corbett1995knowledge,anderson2001revised}. K5: the DAG, its acyclicity guarantee and transitive reduction are what make the gating policy implementable. K8: the design engages directly with the research literature on BKT variants \cite{baker2008more,yudelson2013individualized,pardos2011kt} and on prerequisite-structured student models \cite{kaser2017dynamic}. \\
    P2 & \textbf{Range of conflicting requirements.} Involves wide-ranging or conflicting technical, engineering, and other issues. &
    Granularity conflicts with question-bank coverage: finer skills localise gaps better but multiply the items needed (the rebuilt ontology raised the generation space from $430\times6$ to $1{,}688\times6$ tuples). Diagnostic length conflicts with learner patience: 30 questions cannot cover 430 skills, forcing inference by propagation rather than measurement. Adaptivity conflicts with stability: an aggressive learning rate tracks change quickly but overreacts to a lucky guess. \Cref{sec:alternatives} resolves these explicitly. \\
    P3 & \textbf{Depth of analysis required.} Has no obvious solution and requires abstract thinking and originality in analysis to formulate suitable models. &
    Textbook BKT is per-skill and independent; it cannot express ``mastery of $s_j$ is impossible while $s_i$ is unmastered.'' Making it graph-aware required an original three-phase policy (Bloom-conditioned update, ancestor pull-up, conjunctive transition gating) that is not a published algorithm but a composition designed for this problem (\Cref{sec:mastery-policy}). \\
    P4 & \textbf{Familiarity of issues.} Involves infrequently encountered issues. &
    Bangla-language STEM content with embedded \LaTeX{} mathematics is outside the well-trodden path: it broke JSON escaping (a corruption affecting a measured 13.4\% of one generated batch), broke text rendering in the browser, and is poorly served by English-centric embedding models. \\
    P5 & \textbf{Extent of applicable codes.} Is outside problems encompassed by standards and codes of practice for professional engineering. &
    No standard governs what a ``skill'' is, how finely a syllabus should be decomposed, or when an automatically generated assessment item is fit to serve to a student. These decisions were made against explicit project-defined criteria (\Cref{sec:standards}) because no code of practice supplies them. \\
    P6 & \textbf{Extent of stakeholder involvement and conflicting requirements.} Involves diverse groups of stakeholders with widely varying needs. &
    Learners want the shortest path to a higher score; teachers want syllabus fidelity and correct content; the system's own policy wants maximal information per question. These pull in different directions --- the most informative question is often not the most encouraging one. \Cref{sec:requirements} tabulates the conflicts. \\
    P7 & \textbf{Interdependence.} Is a high-level problem including many component parts or sub-problems. &
    Five interdependent subsystems: the knowledge ontology, the inference engine, the question bank and its generation pipeline, the serving API, and the learner-facing interface. The ontology is upstream of all others --- its granularity determines what the engine can infer and what the generator must produce, which is why a sparse first ontology forced a full rebuild (\Cref{sec:ontology-rebuild}). \\
\end{longtable}}

\begin{table}[htbp]
  \centering
  \caption{Knowledge profiles referenced by P1}
  \label{tab:knowledge-profile}
  \begin{tabularx}{\textwidth}{lY}
    \toprule
    Code & Description \\
    \midrule
    K3 & A systematic, theory-based formulation of engineering fundamentals required in the engineering discipline. \\
    K4 & Engineering specialist knowledge that provides theoretical frameworks and bodies of knowledge for the accepted practice areas in the engineering discipline; much is at the forefront of the discipline. \\
    K5 & Knowledge that supports engineering design in a practice area. \\
    K6 & Knowledge of engineering practice (technology) in the practice areas in the engineering discipline. \\
    K8 & Engagement with selected knowledge in the research literature of the discipline. \\
    \bottomrule
  \end{tabularx}
\end{table}

\section{Project Objectives and Scope}
\label{sec:objectives}

\subsection{Objectives}

The project set the following verifiable objectives. \Cref{ch:evaluation}
evaluates the delivered system against each.

\begin{enumerate}[label=\textbf{O\arabic*.}, leftmargin=*]
  \item \textbf{Construct a prerequisite skill graph} for the HSC Mathematics,
    Physics and Chemistry syllabus at a granularity fine enough that a single
    node names one independently testable action, and verify that the graph is
    acyclic and free of isolated nodes.
  \item \textbf{Implement a knowledge-tracing engine} that maintains a
    per-learner, per-skill mastery estimate, updates it from each response by
    Bayesian inference, and conditions that inference on the cognitive level of
    the question.
  \item \textbf{Make the estimate graph-aware}, such that evidence propagates to
    prerequisite skills and no skill can be credited above an unmastered
    prerequisite.
  \item \textbf{Deliver an adaptive diagnostic} that localises a learner's
    frontier within a bounded number of questions (target 30) rather than
    testing every skill.
  \item \textbf{Deliver targeted practice} that drives a chosen topic to a
    defined mastery threshold and detects when repeated failure is caused by a
    missing prerequisite rather than by the topic itself.
  \item \textbf{Produce a question bank} aligned to the skill graph, in Bangla
    with correct mathematical typesetting, via an automated pipeline with a
    verification stage.
  \item \textbf{Deliver a usable web application} exposing the diagnostic,
    practice and a mastery view to a learner.
  \item \textbf{Verify the engine end to end} with an automated test harness that
    runs without external credentials.
\end{enumerate}

\subsection{Scope}

\textbf{Included:} the three admission-test subjects; multiple-choice items
with four options; a single-process web deployment sufficient for demonstration;
Bangla content with \LaTeX{} mathematics; English and Bangla interface text.

\textbf{Explicitly excluded:} written/subjective answer grading; handwriting or
image input; mobile-native applications; teacher-facing classroom analytics;
multi-tenant institutional deployment; live proctoring; and any claim of
improved examination outcomes, which would require a controlled study with human
participants that this project did not undertake (\Cref{sec:eval-limits}).
